\scanpage{88}
% archive foliation: 87
\deleted{\unsure{(v)|(iv)}}\deleted{\unsure{(iv)|(iii)}}
i.e. la forme induite sur $M$ \inserted{par une forme ample}
est non dégénérée,
donc $N=M\oplus M^\perp$, et la forme sur
$N$ est la somme des formes induites sur
$M$, $M^\perp$.

(v) La somme directe de deux formes
\deleted{positives \unsure{ou}} pseudo-amples (amples) est
itou.

(vi) Le produit tensoriel de deux formes
amples est ample.

(vii) Si $\rho=0$, de sorte que essentiellement $M$
est défini par un vectoriel de dim
finie sur $\mathbf Q$, alors les formes amples
sont les formes \uline{définies positives}.
\deleted{\unsure{Énoncé} analogue pour $M(\rho)$, $M$ un vectoriel
de dim finie sur $\mathbf Q$; utilisant
l'isom canon}

(viii) Cas de $\mathbf Q(\rho)$ : on a $\Phi^s(\mathbf Q(\rho))\simeq\mathbf Q$,)
et on veut que les formes amples
correspondent aux nbs \inserted{\uline{strict}} \uline{positifs}.
